\documentclass{subfiles}
\begin{document}
    We have discussed so far of how the securety of our algorithms is given by,
        among other things, the necessifty of computing big numbers reduced in modulo.
        The question then is, can we ease the task when the user has rights on the data?
        The affirmative answer is provided by the \emph{Square-and-Multiply} algorithm.
        \emph{Figure \ref{Fig:1}} shows the pseudo-code of the algorithm.

    Informaly speaking, given a number in exponential form; 
        we decompose the exponent as subsequent powers of two,
        compute the reduced exponentiations and multiply the results, reducing in modulo 
        accordingly.
        \subfile{../../extra/TikZ/Figure 1 - Pseudo code of SaM}

    More rigorously, let \(g\) be the base number, let \(e\) be the exponent
        and \(n\) the modulo. That is, we want to compute
        \[
            g^{e} \mod{n}.
        \]
        For the sake of our  discussion, assume \(e >> 1\).
    
    Essentially, we consider all the subsequent powers of 2, up to 
        \(2^{\log{e}}\). Write \(e\) in it's binary expansion and 
        compute \(g^{e}\) as the product of \(g^{2^{i}}\), 
        where \(2^{i}\) is in the binary expansion of \(e\).
\end{document}
